\documentclass{article}
\usepackage[english,greek]{babel}
\usepackage[utf8]{inputenc}
\def\tl{\textlatin}
\def\tg{\textgreek}
\title{ΠΑΝΕΠΙΣΤΗΜΙΟ ΤΟΥ \tl{MICHIGAN} ΙΝΣΤΙΤΟΥΤΟ ΑΓΓΛΙΚΗΣ ΓΛΩΣΣΑΣ}
\begin{document}
\date{} % clear date
\author{} %clear author
\maketitle
Αν Αρμπόρ (\tl{Ann Arbor}), Μίσιγκαν – Ηνωμένες Πολιτείες Αμερικής\\
\begin{center}
    Πιστοποιητικό Επάρκειας στην Αγγλική Γλώσσα
\end{center} 
Με το παρόν πιστοποιείται ότι o Ονοματεπώνυμο πέρασε με επιτυχία τις Εξετάσεις για το Πιστοποιητικό Επάρκειας στην Αγγλική Γλώσσα, ως δεύτερη ή ξένη γλώσσα, στις Ημερομηνία έκδοσης πιστοποιητικού, στη Τόπος, Ελλάδα. \\
\begin{center}
\begin{tabular}{ c c }
 [Υπογραφή] & [Υπογραφή]  \\ 
 Πρόεδρος &  Διευθυντής του Ινστιτούτου Αγγλικής γλώσσας \\  
\end{tabular}
\end{center}    
\begin{center}
\begin{tabular}{ c c }
 [Υπογραφή] & [Υπογραφή]  \\ 
Γραμματέας & Πιστοποίηση γνώσεως της αγγλικής γλώσσας\\  
\end{tabular}
\end{center}    
\begin{center}
    [Έμβλημα]
\end{center} 
Το παρόν αποτελεί ακριβή μετάφραση στην ελληνική από την αγγλική γλώσσα του εγγράφου, του οποίου νομίμως επικυρωμένο αντίγραφο επισυνάπτεται στην παρούσα, εκδίδω δε την μετάφραση αυτή, σύμφωνα με το άρθρο 36 παρ. 2 γ΄ του Κώδικα Δικηγόρων (ν. 4194/2013), βεβαιώνοντας συγχρόνως ότι έχω επαρκή γνώση της γλώσσας από και προς την οποία μεταφράζω. Η ως άνω μετάφραση έχει πλήρη ισχύ έναντι οποιασδήποτε δικαστικής ή άλλης Αρχής, σύμφωνα με το άρθρο 36 παρ. 2 γ΄ του άνω Κώδικα. Τόπος, Ημερομηνία επικύρωσης. Ο/η βεβαιών/ούσα και μεταφράσας/σασα Δικηγόρος
\end{document}
